\documentclass[11pt]{article}
\usepackage[a4paper, top=18mm, bottom=18mm, left=20mm, right=20mm]{geometry}
\usepackage[T2A]{fontenc}
\usepackage[utf8]{inputenc}
\usepackage[ukrainian]{babel}
\usepackage{graphicx}
\usepackage[export]{adjustbox}
\usepackage{amsmath}
\usepackage{amssymb}
\graphicspath{{/app/Quiz/Scripts/2023/ProgAll/15BinTree/}}
\setlength{\parindent}{0pt}
\setlength{\parskip}{6pt}
\emergencystretch=3em
\pagestyle{empty}
\begin{document}
%begin
{\centering\Large\bfseries Контрольна робота\par}
\medskip
\noindent\rule{\linewidth}{0.5pt}
\smallskip
\noindent\textbf{Варіант \No\,3}\hfill \textit{ПІБ:}\,\underline{\hspace{60mm}}
\vspace{4mm}

%beginitem
1. \textbf{Що буде виведено за виконання фрагменту коду?}

%code
print('R', end=' ')
m=4
while m<=10:
    m+=2
print(m)
%code

\par\vspace{5mm}
%enditem

%beginitem
2. 
Вкажіть значення та тип виразу:
"6.0j"==6
\par\vspace{5mm}
%enditem

%beginitem
3. \textbf{Що буде виведено за виконання фрагменту коду?}
try:
    res = 8//8/8*8
    if isinstance(res, bool):
        print(f'bool;{res}')
    elif isinstance(res, int):
        print(f'int;{res}')
    elif isinstance(res, float):
        print(f'float;{res:.2f}')
    else:
        print('???')
except:
    print('Z')
\par\vspace{5mm}
%enditem

%beginitem
4. \textbf{Що буде виведено за виконання фрагменту коду?}
try:
    res = 4**-0.5*True
    if isinstance(res, bool):
        print(f'bool;{res}')
    elif isinstance(res, int):
        print(f'int;{res}')
    elif isinstance(res, float):
        print(f'float;{res:.2f}')
    else:
        print('???')
except:
    print('Z')
\par\vspace{5mm}
%enditem

%beginitem
5. \textbf{Що буде виведено за виконання фрагменту коду?}
try:
    for e in range(7, 7, -2):
        if e<=7:
            continue
        print(e, end=';');e = -5
        if e<=7:
            break
    else:
        print(e,end=';')
    print(e)
except:
    print('Z')
\par\vspace{5mm}
%enditem

%beginitem
6. %goodpath:Scripts/2019/Mod2019_1/Lex/03-good.txt
\textbf{Відмітити все, що є коректним оператором С++:}
( 1 ) x=y=z

( 2 ) ++n

( 3 ) True=1

( 4 ) n=1; n++
\par\vspace{5mm}
%enditem

%beginitem
7. %goodpath:Scripts/2019/Mod2019_1/Lex/01-good.txt
\textbf{Відмітити все, що є однією правильною лексемою:}
( 1 ) 0.5

( 2 ) or

( 3 ) -3J

( 4 ) "abc"""
\par\vspace{5mm}
%enditem

%beginitem
8. \textbf{Що буде виведено за виконання фрагменту коду?}

%code
cout << 9 / 9 / 9 / 9;
%code
\par\vspace{5mm}
%enditem

%beginitem
9. %goodpath:Scripts/2018/Mod2018_1/01-good.txt
\textbf{Відмітити все, що є однією правильною лексемою:}
( 1 ) False

( 2 ) fop

( 3 ) <>

( 4 ) '123'
\par\vspace{5mm}
%enditem

%beginitem
10. \textbf{Що буде виведено за виконання фрагменту коду?}
%code
#include <iostream>
using namespace std;

int f(int &x, int &y){
    x = 6;
    y-= 3;
    return y;
}

int main(){
    int a = 1, b = 4;
    b = f(b, b);
    cout << a << ":" << b <<':';
    {
        int a = 7, b = 2;
        cout << ((b<=4) || ((a-=2) <= 4))
             << a << ":" << b << ":";
    }
    {
        inta = 8;
        cout << a << ':';
    }
    cout << a << ":" << b << endl;
    return 0;
}
%code

\par\vspace{5mm}
%enditem

%end
\newpage
\end{document}

